% Driveplexity - JAIIO 2026 short paper (work in development, max 4 pages
% sin contar pagina titulo + abstracts).
% Anonymous submission.
%
\documentclass{llncs}

\usepackage[a4paper, top=5.2cm, bottom=5.2cm, left=4.4cm, right=4.4cm]{geometry}
\usepackage[T1]{fontenc}
\usepackage[utf8]{inputenc}
\usepackage{mathptmx}
\usepackage{graphicx}
\usepackage{float}
\usepackage{amsmath,amssymb}
\usepackage{booktabs}
\usepackage{enumitem}
\usepackage{setspace}

\setlength{\textfloatsep}{8pt plus 2pt minus 2pt}
\setlength{\floatsep}{6pt plus 2pt minus 2pt}
\setlength{\intextsep}{6pt plus 2pt minus 2pt}
\flushbottom

\usepackage[style=apa,maxcitenames=2,mincitenames=1,maxbibnames=2,minbibnames=1]{biblatex}
\DeclareLanguageMapping{spanish}{spanish-apa}
\addbibresource{references.bib}
\setlength{\bibitemsep}{2pt plus 0.5pt}
\renewcommand*{\bibfont}{\scriptsize}

\begin{document}

% ============================================================
% TITLE + ABSTRACTS (pagina separada, no cuenta para 4 paginas)
% ============================================================
\title{Driveplexity: Endogenous Activation via Homeostasis in
LLM-based Multi-Agent Debate\thanks{During the preparation of this
manuscript the authors used GPT-5.4 to improve readability; authors
reviewed and edited the content and assume full responsibility. Full
AI-tooling detail in the project repository.}}

\author{}
\institute{}

\maketitle

\begin{abstract}
LLM-based multi-agent systems tend to lose coherence as interaction
extends ---\emph{context collapse}---, which in adversarial
deliberation can invalidate the debate with social consequences.
Contemporary frameworks (AutoGen, MetaGPT, LangGraph) keep
activation exogenous: when each agent speaks is decided outside its
internal state, a choice that may treat the symptom rather than the
cause. Biology suggests an alternative: internal variables
accumulate tension, trigger a response and return to equilibrium
without external instruction ---homeostasis. We propose
\textbf{Driveplexity} (DPLXY), an endogenous orchestration
strategy where each agent accrues an epistemic deficit $\delta_i$
that drifts upward during inactivity and is reduced after a useful
contribution; a sigmoid turns it into action probability. A linear
TD(0) Q-learner is added as an experimental extension whose
convergence remains an open question. Explicit \emph{Axioms of
Homeostatic Autonomy} (AAH, A1--A3) state its ontological
commitments. We contrast DPLXY with a LangGraph router in
adversarial political debate ($n{=}5$ pairs, GPT-5-nano),
alongside a \emph{DPLXY-no-Q} ablation ($n{=}1$) and a
\emph{LangGraph-informed} fair baseline ($n{=}2$). The data are
compatible with DPLXY reducing floor monopolization (top share
$\sim$1/8 vs $\sim$1/3 under LangGraph, preserved under volume
matching in both early and late windows) and being directionally
associated with higher AAF acceptance ($\Delta\approx{+}0.09$ to
${+}0.18$; preliminary, non-significant), even while the TD loop
fails to converge. The two ablations are compatible with the
effect residing in the regulatory closure rather than in learning
or information access: \emph{DPLXY-no-Q} reproduces or intensifies
the macro pattern, and \emph{LangGraph-informed} ---given the same
internal features--- does not close the gap with DPLXY.

\keywords{multi-agent systems \and large language models \and
homeostatic activation \and endogenous orchestration \and
multi-agent debate}
\end{abstract}

\clearpage

\title{Driveplexity: Activaci\'on End\'ogena por Homeostasis en
Debate Multi-Agente con LLMs}
\maketitle

\begin{abstract}
Los sistemas multiagente basados en LLMs tienden a perder coherencia
al extenderse la interacci\'on ---\emph{colapso del contexto}---,
fen\'omeno que en deliberaci\'on adversarial puede invalidar el
debate con consecuencias sociales. Los frameworks contempor\'aneos
(AutoGen, MetaGPT, LangGraph) mantienen la activaci\'on ex\'ogena:
cu\'ando habla cada agente se decide por fuera de su estado
interno, decisi\'on que podr\'ia tratar el s\'intoma y no la causa.
La biolog\'ia sugiere una alternativa: variables internas acumulan
tensi\'on, disparan respuesta y retornan al equilibrio sin
instrucci\'on externa ---la homeostasis. Proponemos
\textbf{Driveplexity} (DPLXY), una estrategia end\'ogena en la
que cada agente acumula un d\'eficit epist\'emico $\delta_i$ que
deriva hacia arriba al permanecer inactivo y se reduce tras una
contribuci\'on \'util; una sigmoide lo traduce en probabilidad de
acci\'on. Un Q-learner lineal TD(0) se incluye como extensi\'on
experimental cuya convergencia queda como pregunta abierta. Los
\emph{Axiomas de Autonom\'ia Homeost\'atica} (AAH, A1--A3)
explicitan sus compromisos ontol\'ogicos. Contrastamos DPLXY con
un enrutador LangGraph en debate pol\'itico adversarial ($n{=}5$
pares, GPT-5-nano), ablaci\'on \emph{DPLXY-sin-Q} ($n{=}1$) y
l\'inea base \emph{LangGraph-informada} ($n{=}2$). Los datos son
compatibles con que DPLXY reduzca la concentraci\'on de turnos
(cuota m\'axima $\sim$1/8 vs $\sim$1/3 en LangGraph, preservada
bajo volume-matching en ventanas temprana y tard\'ia) y se asocie
direccionalmente con mayor aceptaci\'on AAF
($\Delta\approx{+}0.09$ a ${+}0.18$; preliminar, no significativo),
incluso con el lazo TD sin converger. Las dos ablaciones son
compatibles con que el efecto resida en el cierre regulatorio
antes que en el aprendizaje o en el acceso informacional:
\emph{DPLXY-sin-Q} reproduce o intensifica el patr\'on macro, y
\emph{LangGraph-informado} ---dotado de los mismos rasgos
internos--- no cierra la brecha con DPLXY.

\keywords{sistemas multi-agente \and modelos de lenguaje \and
activaci\'on homeost\'atica \and orquestaci\'on end\'ogena \and
debate multi-agente}
\end{abstract}

\clearpage
\setcounter{page}{1}

% ============================================================
% Pag 1: motivacion + research question
% ============================================================
\section{Motivaci\'on}
\label{sec:motivation}

Los sistemas multiagente basados en LLMs prometen sostener
deliberaci\'on extendida en pol\'iticas p\'ublicas, revisi\'on
cient\'ifica o soporte institucional. Dominios en los que un debate
viciado afecta decisiones que trascienden al software. En horizontes
prolongados, de hecho, suelen reportarse p\'erdidas de registro
argumentativo, repetici\'on de puntos y mon\'ologos
desconectados~\parencite{park2023generative}: el \emph{colapso del
contexto}~\parencite{huang2024masresilience}.

El paradigma dominante lo aborda por fuera del agente. Los
\emph{frameworks} contempor\'aneos ---AutoGen~\parencite{wu2023autogen},
MetaGPT~\parencite{hong2023metagpt}, LangGraph--- estructuran la
interacci\'on con grafos y enrutadores centralizados, y el
\emph{heartbeat}~\parencite{openclaw2025heartbeat} suma persistencia
temporal; en todos, \emph{cu\'ando} habla cada agente se decide fuera
de su estado interno. Esa activaci\'on ex\'ogena no captura la
presi\'on epist\'emica que cada agente acumula y podr\'ia tratar el
s\'intoma sin explicar la causa. Un mecanismo bien establecido en
biolog\'ia describe c\'omo los sistemas regulan desequilibrios sin
control externo ---la homeostasis, formalizada como reducci\'on de
impulso por HRRL~\parencite{keramati2014homeostatic}. A primera vista,
un mecanismo visceral parece ajeno a orquestar agentes de software;
la analog\'ia estructural, sin embargo, es precisa: donde un
organismo acumula hambre y act\'ua para reducirla, un agente
deliberativo podr\'ia acumular d\'eficit epist\'emico y actuar para
reducirlo. De ah\'i la \emph{l\'inea de investigaci\'on}: ¿puede un
agente con estado regulatorio interno sostener acci\'on valiosa y
coherente sin que un planificador externo decida cu\'ando participa?
La \emph{l\'inea de desarrollo} articula una teor\'ia ontol\'ogica,
una arquitectura, un protocolo experimental y \emph{resultados
preliminares} sobre cinco semillas, con \emph{acciones futuras} que
derivan de ellos.

% ============================================================
% Pag 2: TAH compacta + arquitectura
% ============================================================
\section{Axiomas de Autonom\'ia Homeost\'atica}
\label{sec:tah}

Responder la pregunta exige precisar \emph{qu\'e constituye un
agente aut\'onomo}, \emph{qu\'e lo hace actuar} y \emph{qu\'e hace
valiosa una acci\'on} ---supuestos que MAS-LLM suele dejar
impl\'icitos. A1--A3 no son un teorema: formalizan los supuestos
m\'inimos sobre los que se compromete el dise\~no.

\textsc{A1, Autonom\'ia}~\parencite{ashby1956introduction,beer1972brain}
---\emph{¿qu\'e es autonom\'ia?}---: que la pol\'itica del agente
dependa de su estado interno.
\vspace{-0.6em}
\begin{equation}
\mathrm{Aut}(a_i) \iff \exists\,\delta_i \in S_L(a_i):\;
\pi_i = \pi_i\!\left(\delta_i(t)\right). \label{eq:a1}
\end{equation}
\vspace{-1.4em}

\textsc{A2, Impulso}~\parencite{sterling1988allostasis}
---\emph{¿qu\'e hace a un agente actuar?}---: un impulso interno que
crece con el desv\'io del equilibrio y cuya magnitud determina
mon\'otonamente la probabilidad de activaci\'on.
\vspace{-0.6em}
\begin{equation}
D(\varepsilon){=}0,\quad D'(\delta_i){>}0\ \forall\delta_i{>}\varepsilon,\quad
p_i(t) = \sigma\!\left(D(\delta_i(t))\right). \label{eq:a2}
\end{equation}
\vspace{-1.4em}

\textsc{A3, Compuerta de calidad}~\parencite{ryan2001happiness}
---\emph{¿qu\'e hace valiosa a una acci\'on?}---: un cambio efectivo
$g(e^{(t)},e^{(t+1)}){\ge}0$ en el estado del entorno regulado; el
impulso decrece proporcionalmente a ese cambio.
\vspace{-0.6em}
\begin{equation}
\Delta\delta_i = -\alpha\cdot g\!\left(e^{(t)},e^{(t+1)}\right),\quad
\alpha{>}0. \label{eq:a3}
\end{equation}
\vspace{-1.4em}

\textsc{Pr1} (derivada de A1--A3). Dado un agente expuesto a
perturbaciones, cuyo impulso crece con el desv\'io del equilibrio y
s\'olo decrece ante cambio efectivo, se esperar\'ia que emerja
acci\'on sostenida de calidad, con probabilidad de activaci\'on
creciente en el d\'eficit. La implementaci\'on
(\S\ref{sec:agent-arch}) usa aproximaci\'on lineal off-policy ---la
\emph{tr\'iada mortal}~\parencite{sutton2018reinforcement}--- por lo
que la convergencia TD queda como pregunta emp\'irica, no como
supuesto.

\subsection{Arquitectura propuesta del agente}
\label{sec:agent-arch}

DPLXY operacionaliza los AAH en un MAS-LLM. A2 instancia un
d\'eficit epist\'emico
$\delta_i{\in}[\varepsilon,\infty)$ que deriva al permanecer
inactivo y decae por saciedad
$\delta{\leftarrow}\max(\varepsilon,\delta{-}\alpha\Delta\varphi^*)$;
una sigmoide lo traduce en probabilidad de acci\'on y la
recompensa realiza A3 con $D(\delta){=}(\delta{-}\varepsilon)^2$.
Un \emph{StimulusEvaluator} determinista fusiona cinco rasgos
estructurales ---relevancia faccional, centralidad del nodo
atacado, novedad, presi\'on no respondida y memoria sem\'antica---
por suma ponderada con pesos a priori, modulando $\delta_i$. La
calidad $\Delta\varphi^*$ se define sobre un
AAF~\parencite{dung1995acceptability} (centralidad, ciclos
dial\'ecticos, diversidad); es insumo de saciedad y observable
diagn\'ostico, circularidad parcial que obliga a priorizar
m\'etricas estructuralmente independientes
(\S\ref{sec:resultados}). Entre acciones (\textsc{debate},
\textsc{search}, \textsc{read}, \textsc{message}, \textsc{pass})
se selecciona por Q-learning lineal TD(0) sobre estado
normalizado; el aprendizaje se incluye como extensi\'on
experimental no cr\'itica al cierre regulatorio.

% ============================================================
% Pag 3: diseno + resultados preliminares
% ============================================================
\section{Dise\~no experimental}

La arena debe forzar interdependencia argumentativa ---para que la
p\'erdida de coherencia sea observable--- y mantenerse acotada en
c\'omputo. Adoptamos debate pol\'itico de suma cero inspirado en
Sotopia~\parencite{zhou2024sotopia} durante $T{=}80$
\emph{heartbeats} entre diez agentes
GPT-5-nano~\parencite{openai2025gpt5nano} en dos facciones espejo
(Gobierno/Oposici\'on), cada una mapeada a los cinco subsistemas
del VSM~\parencite{beer1972brain} (S1 operaci\'on, S2 coordinaci\'on,
S3 control, S4 inteligencia, S5 pol\'itica) como unidad m\'inima
de organizaci\'on social, lo que habilita buscar patrones
emergentes por rol adem\'as de la regulaci\'on individual.

La variable experimental es la estrategia de activaci\'on:
\textbf{DPLXY} (end\'ogena: sigmoide homeost\'atica) vs
\textbf{LangGraph} (ex\'ogena: enrutador LLM neutral). Ese contraste
admite dos lecturas alternativas que el resto del dise\~no busca
cerrar. La lectura \emph{informacional} ---DPLXY accede a rasgos
y $\delta_i$ que LangGraph no ve--- se ataca con
\textbf{LangGraph-informado}, que recibe los mismos rasgos por
JSON; la lectura \emph{mecan\'istica} ---el efecto podr\'ia
provenir s\'olo del cierre regulatorio determinista--- se ataca
con \textbf{DPLXY-sin-Q}, que desactiva el aprendizaje y conserva
la regulaci\'on. Cinco semillas $\{7,17,99,123,256\}$ pareadas
por Wilcoxon \emph{signed-rank} alimentan el contraste base; los
controles se corren sobre un subconjunto por costo. Los
hiperpar\'ametros se fijan \emph{antes} sobre la semilla 42 con
dos criterios a priori ---densidad en rango operativo y
estabilidad del d\'eficit alrededor del \emph{set-point}
$\varepsilon$--- y se congelan; la semilla 42 queda \emph{excluida}
de la evaluaci\'on para prevenir fuga de calibraci\'on. Igualar
tokens o aplicar \emph{rate limit} a DPLXY anular\'ia la
auto-activaci\'on que se quiere medir; por eso adoptamos
\emph{heartbeats id\'enticos} y tratamos la asimetr\'ia de volumen
con control \emph{volume-matched} en dos ventanas
(\S\ref{sec:resultados}).

\section{Resultados preliminares}
\label{sec:resultados}

Dos m\'etricas estructuralmente independientes del sensor sostienen
el contraste: \emph{aceptaci\'on AAF} (sem\'antica
grounded~\parencite{dung1995acceptability}, proxy de resiliencia)
y \emph{cuota m\'axima} (proxy de monopolizaci\'on).
$\mathrm{PRR}_G$ satura $>\!0.98$ en todas las condiciones (control
nulo); $\Delta\varphi^*$ queda como auxiliar por circularidad
parcial con el sensor. Una evaluaci\'on externa (juez humano o
LLM-as-judge) queda como acci\'on futura. La
Tabla~\ref{tab:results-short} sintetiza; cada bloque contrasta una
observaci\'on contra lecturas alternativas.

\begin{table}[H]
\caption{Resultados preliminares. DPLXY/LangGraph: $n{=}5$ semillas
$\{7,17,99,123,256\}$; DPLXY-sin-Q: $n{=}1$; LangGraph-informado:
$n{=}2$ (seed 42 excluida por fuga de calibraci\'on, $T{=}80$).
\emph{vm-1}: primeros $K{=}N_{\text{LG}}$ debates de DPLXY;
\emph{vm-$\ell$}: \'ultimos $K$. Ablaciones en m\'etricas
\emph{full} ($\dagger$). M\'etricas vm computadas sobre
$\{7,17,123,256\}$ hasta cierre de la corrida 99. Media $\pm$
desviaci\'on muestral.}
\label{tab:results-short}
\centering\small
\begin{tabular}{|l|c|c|c|c|}
\hline
\textbf{M\'etrica} & \textbf{DPLXY} & \textbf{LangGraph} & \textbf{DPLXY$-$Q} & \textbf{LG-inf.} \\
\hline
Total debates                      & $384\pm6$        & $72\pm6$          & $604$            & $75\pm0$         \\
Densidad (deb./hb.)                & $4.80$           & $0.90$            & $7.55$           & $0.94$           \\
Cuota m\'ax.\ (vm-1)               & $0.133\pm0.007$  & $0.349\pm0.052$   & $0.106^\dagger$  & $0.260\pm0.010^\dagger$ \\
Cuota m\'ax.\ (vm-$\ell$)          & $0.158\pm0.009$  & $0.349\pm0.052$   & --               & --               \\
AAF acc.\ (vm-1)                   & $0.640\pm0.035$  & $0.546\pm0.023$   & $0.729^\dagger$  & $0.540\pm0.010^\dagger$ \\
AAF acc.\ (vm-$\ell$)              & $0.724\pm0.055$  & $0.546\pm0.023$   & --               & --               \\
\hline
\end{tabular}
\end{table}

\paragraph{¿Se sostiene la acci\'on?} DPLXY produce $4.8$
debates/heartbeat frente a $0.9$ en LangGraph ---factor
$5.4\times$ a igual presupuesto temporal, no a igual presupuesto de
tokens, ya que el enrutador consume los suyos y la comparaci\'on se
lee por heartbeat. La presi\'on end\'ogena sugiere una tendencia a
sostener la actividad a lo largo del horizonte, all\'i donde la
orquestaci\'on ex\'ogena decae. Queda por determinar si esa
actividad sostiene coherencia o s\'olo suma ruido.

\paragraph{¿Sostiene coherencia bajo control de volumen?} Truncando
DPLXY a $K{=}N_{\text{LG}}$ debates por semilla, \emph{primeros-$K$}
(contexto corto) da $0.640$ vs $0.546$ ($\Delta{\approx}{+}0.09$);
\emph{\'ultimos-$K$} (contexto equivalente a LangGraph) sube a
$0.724$ ($\Delta{\approx}{+}0.18$). La direcci\'on se sostiene,
aunque vm-$\ell$ no neutraliza el \emph{burn-in}: los \'ultimos-$K$
de DPLXY ocurren tras ${\sim}300$ debates previos ---una
condici\'on sim\'etrica con presupuesto extendido para LangGraph
queda como acci\'on futura. Wilcoxon pareado ($W{=}0$,
$p{=}0.0625$) es piso mec\'anico para $n{=}5$: no cruza
$\alpha_{\text{adj}}{=}0.025$ con independencia del efecto.

\paragraph{¿C\'omo se sostiene? Reparto de turnos.} Podr\'ia
pensarse que DPLXY mejora coherencia silenciando a los d\'ebiles,
o que su cuota baja sea artefacto del mayor volumen. La
\emph{cuota m\'axima} es inconsistente con ambas: macro cae de
$0.362$ (LG) a $0.126$ ---cerca del uniforme $1/N{=}0.10$, con
$\sigma$ un orden de magnitud m\'as estable ($0.006$ vs
$0.052$)--- y se mantiene baja en ambas ventanas vm (vm-1:
$0.133$; vm-$\ell$: $0.158$). Los datos son consistentes con una
din\'amica end\'ogena donde cada agente contribuye cuando su
$\delta_i$ lo justifica, distinta del r\'egimen ex\'ogeno
``quien ya habl\'o vuelve a hablar''.

\paragraph{¿Qu\'e componente sostiene el efecto? Ablaciones.} La
recompensa TD media es negativa ($-2.03\pm0.38$): el lazo Q no
converge, compatible con la tr\'iada mortal. Las ablaciones
---$n{=}1$ y $n{=}2$, evidencia anecd\'otica--- acotan
direccionalmente las lecturas alternativas. \emph{DPLXY-sin-Q}
reproduce o intensifica el patr\'on macro (cuota $0.106$, AAF
$0.729$): el aprendizaje no aparece como condici\'on necesaria.
\emph{LangGraph-informado}, con acceso a los mismos rasgos y a
$\delta_i$ por JSON, mantiene cuota $0.260\pm0.010$ y AAF
$0.540\pm0.010$ ---indistinguibles de LangGraph base---: entregar
la informaci\'on no acerca al r\'egimen end\'ogeno. La lectura
parsimoniosa es compatible con que el efecto resida en el cierre
regulatorio (A2{+}A3) antes que en el aprendizaje o el acceso
informacional; confirmarlo exige escalar los controles.

% ============================================================
% Pag 4: discusion + future work
% ============================================================
\section{Discusi\'on, limitaciones y acciones futuras}
\label{sec:discussion}

\paragraph{Alcance y acciones futuras.} Con $n{=}5$ base y
$n{=}1,2$ en ablaciones, Wilcoxon no cruza Bonferroni aun con
direcci\'on consistente: reportamos tendencia, no prueba. El dual
primeros-$K$/\'ultimos-$K$ acota el confound de contexto y las
ablaciones los de mecanismo e informaci\'on; sensibilidad a
$(\theta,k)$, ablaci\'on VSM y burn-in de vm-$\ell$ quedan
abiertos. La evaluaci\'on se apoya en un \'unico LLM y dominio
adversarial de suma cero. Como \emph{acciones futuras}: escalar a
$n{\geq}20$ con $k$-fold; sustituir \emph{heartbeats id\'enticos}
por \emph{rate limit} que iguale densidad; estudiar patrones
sociales emergentes por rol VSM; y replicar en dominios
cooperativos y multi-LLM. Si la l\'inea resiste mayor escrutinio,
parte del esfuerzo actual en orquestaci\'on ex\'ogena podr\'ia
estar tratando un s\'intoma antes que la causa. C\'odigo, semillas
y logs se liberar\'an al desanonimizar.

\printbibliography

\end{document}
